% The introduction. Written after the notes pass of 2026-08, out of the findings
% recorded in apparatus/print-notes.md; every page cited here has a note against it.
% ⚠ It must not restate the PREFACE, which carries the editorial matter: the gap in
% print, the textual policy, the Latin question, the English register and Coverdale,
% the page, the unprinted Hebrew, Brightman. This piece is a reading of the book.
% If you change a claim here, change it in the note it came from.
\begingroup
\setlength{\parskip}{0.55\baselineskip}
\setlength{\parindent}{0pt}
\addcontentsline{toc}{section}{Introduction}
{\centering\Large\scshape Introduction\par}
\vspace{1.5\baselineskip}

This book has been read for three centuries in pieces. It is a manual of private
devotion, and a manual is used by opening it where you need it; the editions have
encouraged that, and the one nineteenth-century translation most readers know
rearranged the whole for exactly that purpose. What follows is an argument for
reading it the other way --- through, in its own order --- and an account of what
that turns up.

Page numbers in what follows are the 1853's, the bracketed figures on the inner
edge of every leaf of this edition.

\vspace{0.8\baselineskip}
{\scshape One thing to notice first.}\quad
The book does not point at itself.

A phrase will be used on page 22 and used again on page 431, turned round the
second time, and neither page marks the other. The night at the ford is quarried
three times, for three different things, and no leaf mentions the other two.
Hezekiah's psalm carries the whole of a confession in Part II, returns in Part III
in a form that differs from it twice, and comes back a third time at the evening
office as a thanksgiving; the three do not acknowledge one another. The body is
catalogued member by member at 406 as filth and again at 426 as workmanship, in
what is plainly one design, and neither page looks up. The last sentence of the
book was already said, four hundred leaves earlier, among the morning prayers.

None of this is marked, and there is no reason it should have been. Andrewes was
not making a book for a reader. He was making one for himself, and a man does not
cross-reference his own memory. But the consequence for us is that almost
everything this volume is doing is invisible to anyone consulting it, and visible
only to someone going through it. It is why the notes at the foot of the
right-hand pages exist, and why they so often say that nothing on either page
marks the return.

\vspace{0.8\baselineskip}
{\scshape The book's own account of itself.}\quad
It does explain itself once, and the page is easy to walk past.

At 347 stands a leaf headed \emph{Divisio}. It is not a prayer. It is an analysis
of what a prayer contains, set out in two columns, and most of its terms are the
actual section headings of this volume: \emph{Confessio peccatorum}, \emph{Confessio
fidei}, \emph{Confessio laudis}, \emph{Deprecatio}, \emph{Intercessio},
\emph{Supplicatio}, \emph{Gratiarum actio}, \emph{Commendatio}. The page is a table
of contents for a book that never prints one. Three of its terms ---
\emph{Compellatio}, \emph{Comprecatio}, \emph{Threni} --- occur nowhere else.

Set that beside the opening. The first three openings of Part I take the
\emph{when}, the \emph{where} and the \emph{how} of prayer in turn, and each is
built by the same method: a number is fixed from Scripture, and then every item
under it is given a text in which somebody is found doing that thing. \emph{Seven
times a day do I praise thee} sets the number of the hours, and each of the seven
gets a text --- Mark's \emph{a great while before day}, the third hour at
Pentecost, Peter's sixth hour on the housetop, the ninth at the Beautiful Gate,
Isaac at eventide. Then the places, by the same rule: the closet, the housetop,
the temple, the sea-shore, the garden, the bed, the wilderness --- and then
\emph{every where}, which is not an eighth place but the conclusion drawn from the
seven. Then the postures, seven of them, each with the text of somebody in that
posture and the affection it belongs to named in a column alongside.

The canonical hours are being \emph{derived}, not inherited. A man who wanted the
offices could have taken them from a breviary; this one builds them out of
Scripture and shows his working. Whatever else the opening of this book is, it is
not a preface. It is a treatise, and it is arguing that the shape of the day's
prayer can be got from the Bible alone.

\vspace{0.8\baselineskip}
{\scshape How it quarries.}\quad
What the book does with a text is never simply to quote it. Four habits recur, and
once they are seen they are everywhere.

\emph{It changes the mood.} A promise is made into a petition. God says, through
Isaiah, \emph{I have blotted out thy transgressions as a cloud}; the prayer answers
\emph{blot out} --- perfect turned imperative, and \emph{thy} turned to \emph{mine}
(22). At 431 the same verse comes back fitted to the hour: blot out our iniquities
as the cloud of night, our sins as the morning mist. At 332 it stands untouched in
God's own indicative. Three settings, and no page marks the other two. At 349 an
indicative is made a question instead: John's \emph{having loved his own, he loved
them unto the end} becomes \emph{For that thou hast loved me hitherto?} At 355 a
promise is not turned but \emph{taken up with its half discharged} --- \emph{open
thy mouth wide, and I will fill it} is answered \emph{Dilato os meum, Domine: Tu
imple}, I have opened it wide: do thou fill it. The condition is reported as
performed and the promise called in.

\emph{It turns a text round on its speaker.} The sharpest instance is §37's
confession, which is built on Isaiah xxxviii, Hezekiah's psalm of the sick king,
quarried out of order --- verses 15, 17, 16, 17, 14 across two leaves --- and with
the person changed. The Vulgate has \emph{quid respondebit mihi, cum ipse fecerit},
what shall he answer me, seeing He hath done it. The page has \emph{cum ipse
fecerim, fecerim, fecerim}: seeing I myself have done it, have done it, have done
it. One letter changes the person, and the repetition is the confession (372).
Part III has the twin at 409 and it differs twice --- \emph{hæc fecerim}, once, and
under a chapter figure one numeral short. The same habit puts \emph{thy will be
done} into the mouth of Balaam, who wanted to do otherwise and could not, and into
Eli's, who has just been told his house is ruined (351); and it takes a word spoken
by David to King Saul --- \emph{wilt thou chase a flea?} --- and speaks it to God
(412).

\emph{It quarries its own vocabulary from the verse it is standing beside.} The six
titles at 340, the four cells at 341, the nine epithets of \emph{fides} at 342 are
each lifted out of the Vulgate of the text printed alongside. Once the habit is
established it becomes evidence: at 341 it convicts a printed reference of a wrong
numeral, because the word on the page could only have come from the other verse.

\emph{And it does the same thing several times without repeating itself.} The
Lord's Prayer is treated three times in Part II and the three are not to be
conformed. §25 sets it out plainly (350). §26 says it seven times over, petition by
petition, in other men's words, and then begins the seven again with different
speakers, and the later sets are each quarried from one region of Scripture --- the
Psalms, then Proverbs, then the Prophets; the petitions are never once restated in
Andrewes' own words (351--354). §31 states each petition, admits it to have been
failed, and then asks it: \emph{to be held holy by all men, and by me first of all
above many} --- \emph{Yet I have not so held it.} The petition becomes an
indictment before it becomes a request (360). Only the third turns the prayer
against the man praying it.

The fourth petition is the one to watch through all of this. It is never simply
asked. It is asked by Jacob's vow at Bethel with the stone still under his head; by
Job's self-curse negated, \emph{let thistles} \emph{not} \emph{grow instead of
wheat}, so that the oath a man called down on himself if he lied becomes the prayer
for bread; and by \emph{not by bread alone}, which grants the petition and takes it
back in one line.

\vspace{0.8\baselineskip}
{\scshape The book argues.}\quad
Part III is where the method becomes a case, and it is the part of the book most
damaged by being read in pieces, because it is continuous for twenty leaves.

The section is a plea for mercy, and it is built like one. Texts are entered as
evidence: the law, the prophets, nature, the apostles, and last of all Christ. The
apostles are called in a numbered series --- \emph{thine Apostle}, then \emph{the
second Apostle}, then \emph{the fourth} --- with James standing between them,
quoted in his place and never counted (423). Then the last witness, and the tense
is the point: \emph{nor in vain shalt Thou Thyself have said, Come unto Me, all ye
that labour} --- so that the whole chain is set as evidence of something that would
be void if the invitation failed. Then the verdict clause: \emph{these things were
not said in vain, they cannot be}. Then three \emph{wherefores}, marking the turn
from proving to asking. And then, after twenty leaves of argument, the shortest
prayer in the Gospels: \emph{God be merciful to me a sinner}. The disproportion is
the design. The whole apparatus was built to earn the right to say the one sentence
that needed none of it (424).

Two features of the arguing are worth naming, because they are unusual and they are
consistent.

The first is that the argument is strongest where the example is worst. Three human
kings are entered as precedents binding upon God, and the series mounts downwards:
David sparing Shimei, and \emph{David was a man after Thine own heart}; then the
kings of Israel, who \emph{are merciful} --- which is not Israel's estimate of them
but Ben-hadad's servants', a pagan's reckoning quoted with approval; and last Ahab,
\emph{who had sold himself unto sin}, whom God spared when he humbled himself. The
final plea is that God should not do less for this man than He once did for Ahab
(412). The same move again at 422, with three men who would not take back a word:
Isaac, who would not revoke a blessing got by fraud; Darius, bound by the law of
the Medes against his own friend; and \emph{profane Pilate}, \emph{what I have
written I have written} --- \emph{surely God never will.} The series descends in
dignity while it ascends in force, and ends on the man who condemned Christ.

The second is that God's goodness is sorted into classes in order to claim the
lowest. \emph{Good to the good and to them that deserve well; gracious to strangers
and to them that deserve nothing; merciful to the evil, and to them that deserve
ill.} Then five words: \emph{Upon this last I take my stand} (415). The taxonomy
exists to put the man praying into the worst bin and hold him there --- and the
ground offered is not desert but its absence: \emph{when for no man's sake Thou
canst, when for nothing's sake, for Thine own sake Thou forgivest sins.}

The same shape governs §28's ladder in Part II, eleven pairs of \emph{si non
\textellipsis\ tamen}, every one a bargain struck downwards. A scriptural ceiling
is set, conceded to be out of reach, and the least he will do is named instead:
seven times a day like David, or at any rate three like Daniel; Solomon's long
prayer, or at any rate the Publican's short one; Christ's whole night, or at any
rate one hour. The find is that the floor is repeatedly set by a bad example: if
not on the ground in ashes, yet not in beds; if not in sackcloth, yet not in purple
and fine linen --- which is the rich man of the parable. Where the ceiling is a
saint the floor is a sinner. If I cannot be Job, let me at least not be Dives.

\vspace{0.8\baselineskip}
{\scshape What rite this came out of.}\quad
A reader who wants to know what kind of churchman kept this book will find the
evidence pointing hard in two directions at once, sometimes within a single
opening. The honest thing is to set both out and adjudicate neither, and that is
what this edition does.

Within eight leaves in Part II the Latin pages carry five independent marks of the
Eastern rite. At 349 the private prayer takes the shape of the eucharistic
\emph{anaphora}: a spiritual sacrifice is offered, received upon a spiritual altar,
and the Spirit asked to be sent down in return, which is an epiclesis --- a man
alone in his chamber praying the shape of the Church's central action over himself.
On the same leaf stands \emph{voluntary and involuntary}, ἑκουσίων καὶ ἀκουσίων,
which is the Greek liturgies' formula and which no English rite has. At 350 the
Lord's Prayer carries its doxology, absent from the Vulgate at Matthew vi. 13 and
from the Roman rite. At 355 Christ is called \emph{Carbo duplicis naturæ}, the coal
of the twofold nature --- a title that is not in Isaiah: the prophet's coal, taken
from the altar with tongs, read as wood and fire at once and so God and man at
once, with the tongs as what handles what cannot be touched. That is the ἄνθραξ of
the Eastern liturgies, where the same coal is what the communicant receives, and it
carries a whole Christology in one noun. And at 356 he does not merely cite the
Greek Psalter but prays in it: \emph{Domine, dilexi decorem domus Tuæ} stands at
the door and returns fourteen lines later as Κύριε, ἠγάπησα εὐπρέπειαν οἴκου σου.

Two leaves later the evidence turns round. At 358 the section headed \emph{Actus
Adorationis} opens with the Litany --- \emph{O God the Father, of heaven; O God the
Son, Redeemer of the world; O God the Holy Ghost; O holy, blessed Trinity} --- the
Prayer Book's first four invocations in the Prayer Book's own order. The Litany's
un-English \emph{of heaven} is a rendering of the very \emph{de cœlis} standing on
this page; both are drawing on the same Latin. At 359 a prayer for the departed is
made in the Requiem's own words, \emph{defunctis requiem et lucem perpetuam}, in
the private manual of a bishop of the Church of England, and the page lays its
ground first: Christ is Lord alike of the living and of the dead, \emph{whose we
are} and \emph{whose also are they}, and the petition divides its gifts along that
line. And at 360 both columns are the medieval memory-verses of the seven works of
mercy, \emph{visito, poto, cibo, redimo, tego, colligo, condo}, one verb to a work
and made to be counted on the fingers --- not composed here but quoted, as every
schoolman knew them.

The single most loaded line in Part II is at 363, and it does not resolve the
question so much as hold it open on purpose. The \emph{Interpellatio Eucharistica}
is seven verbs long: \emph{Recolo, gratias ago, admoneo, recordor, commemoro,
offero, vel peto ut offeras}. \emph{Admoneo} --- I put thee in mind --- points the
memorial at God, which is the patristic understanding of the anamnesis and not the
ordinary Western one. And \emph{offero} is qualified in the same breath by \emph{vel
peto ut offeras}: the sacrifice is offered, and then in a single word the offerer
is left an open question. A whole controversy hangs on \emph{vel}. The English here
keeps it and settles nothing, which is the only defensible thing to do with it.

An account of this book that argues only the Greek is answerable for 358 and 360.
An account that argues only the Prayer Book is answerable for 349 and 355. Both
sets of evidence are in this volume, on facing leaves, and the reader now has them.

\vspace{0.8\baselineskip}
{\scshape What was on the desk.}\quad
A related question can be settled further than one might expect, because the book
convicts itself.

Which Psalter Andrewes had by him has usually been argued from wording, which can
be misremembered. This volume settles it from numerals, which have to be looked up.
At 433 two references answer only on the Coverdale Psalter and not on the
Authorised Version --- \emph{the Lord hath granted His loving-kindness in the
day-time} cited xlii. 10, and \emph{the outgoings of the morning and of the
evening} cited lxv. 8. Somebody had the Prayer Book physically open.

But not only that Psalter, and the proof is inside a single text. One catalogue is
printed twice in this book, in two recensions, and the two cite the same verse by
two different numbers: Part I gives \emph{Psal.} xxii. 17 against the pierced hands
and feet, and 312 gives xxii. 16 against the same words. Seventeen is the Prayer
Book; sixteen is the Authorised Version. At 295 the plate numbers \emph{a wind that
passeth away} as lxxviii. 39, which is the Authorised Version against the Prayer
Book's 40. At 274 \emph{Ps.} cxvii. 25 is impossible on either and can only be the
Vulgate's numbering. Three systems are on show, and the 1853 does not hold to one.

Nor does he stop at three. At 383 a section on praise ends in Hebrew, because no
Latin or English Psalter will say what the passage needs: \RL{לך דומיה תהלה} ---
Psalm lxv. 2 in the Hebrew numbering --- \emph{to thee silence is praise}. The
Vulgate takes the same consonants from another root and gives \emph{Te decet
hymnus}, praise befitteth thee; the Prayer Book has \emph{Thou, O God, art praised
in Sion}. Either would destroy the entry, which exists in order to say that praise
falls silent. So the Hebrew stands, and our English follows the Hebrew.

The same industry shows in the Passion. Printed 309--314 is a four-Gospel harmony
--- Matthew, John, Luke and Mark walked forward in step, each leaf drawing on three
or four evangelists at once. He is not following one and supplementing him. Much of
that density is invisible on the plate, because the continuations are marked only
by \emph{Vers.} and \emph{Ibid.}, and it is one of the places where this edition's
apparatus can show what the printed page cannot.

\vspace{0.8\baselineskip}
{\scshape The end of it.}\quad
The last section returns to the first argument, and does not restate it.

\emph{At evening, and morning, and at noonday will I pray} stands at 431 --- the
psalm the volume opened on at page 2, where the hours were being built out of
Scripture, text by text, and shown to be derivable. Four hundred pages later the
hours are simply kept. The case has been made and is being lived on, which is what
that looks like.

Then the day is walked down. Psalm xci is quarried across two leaves to cover the
whole of it: the terror by night and the thing that walketh about in darkness on
one, the arrow that flieth by day and the noonday devil on the next (432).
Ephesians' rule for a quarrel is turned round and prayed --- \emph{let not the sun
go down upon Thy wrath} --- at the hour when the sun going down is not a figure but
the time the prayer is being said (434). The examination of conscience is carried
on past consciousness: \emph{let my sleep be, as a rest from labour, so a rest from
sin}, and then further, \emph{let me think nothing, even in my dreams, that may
offend Thee} (435).

And then one line is said twice with a single word changed. \emph{By the lamp
burning, to see Thee. By the lamp quenched, to see Thee.} The lamp put out is the
bedside candle, and the next lines say plainly that it is also the life --- the
long sleep, the sleep of death, the bed of the sepulchre, the litter of worms, the
covering of dust. Four names for the grave, and every one of them is a bed; the
figure begun at the pillow is carried the whole way down and never exchanged for
another. The petition is identical on both sides of the line because the request
does not change when the light does.

The book's last words are \emph{into Thy hands, O Lord, I commend my spirit, for
Thou hast redeemed me}. That verse already stood at page 26, among the morning
prayers. The volume opens its day and closes its life with one sentence, and the
lamp two lines above has just made the two the same act.

Nothing on either page marks the return. It is the last and the best instance of
the habit this introduction began with, and it is the reason for reading the book
through.

\vspace{0.8\baselineskip}
{\scshape One consolation, withheld and then given.}\quad
A word on the hardest stretch, for a reader who may need it.

Six leaves of §37 are a confession by a man who cannot repent properly and knows
it. The section is made almost entirely of allusion the plate never marks ---
Egypt's leeks, the prodigal's husks, Paul's \emph{who hath bewitched you}, Jacob at
the ford, Romans vii. 18 word for word --- and its complaint is dryness:
\emph{ariditas}, and the fountain of tears that will not come. He spends the whole
of it unable to weep.

The answer arrives at 383 and it is Augustine, and it refuses the devotional
reading of Romans viii. 26 outright. \emph{Is thy spirit, or mine, unutterable ---
which is oftentimes no spirit at all, oftentimes cold?} The groaning that cannot be
uttered is not the individual's. It is the whole Church's, pooled --- \emph{because
they all make up one Dove}, one beseeching fervently and another lukewarmly,
\emph{hence come the unutterable groanings, out of the groanings of all in common}.

After six leaves of a man failing to produce the thing he thinks is required, the
book tells him it was never his to produce. That is the consolation this volume had
been withholding, and it is placed exactly where it does the most work.

\endgroup
\clearpage
